%FlowSemantics.tex

\ksec{Flows (KerML 8.4.4.10.2)}\label{sec:flows}

\pn  A \texttt{Flow} is both a \texttt{Step} and a \texttt{Connector}.  A \texttt{Flow} is a feature typed by a \texttt{Transfer}.
A \texttt{Flow} with \texttt{ownedEndFeatures} must specialize \texttt{Transfers::flowTransfers}.
\ecaption{}
\begin{lstlisting}
    abstract flow flowTransfers: FlowTransfer[0..*] nonunique subsets transfers {
        doc
        /*
         * flowTransfers is a specialization of transfers restricted to type FlowTransfers.
         * It is the default subsetting for non-succession flows.
         */

        end feature source: Occurrence redefines FlowTransfer::source, transfers::source;
        end feature target: Occurrence redefines FlowTransfer::target, transfers::target;
    }
\end{lstlisting}

\pn \texttt{Transfer.isInstant} should default to true.

\comment{Source and target should be directed features of a  \texttt{Behavior}-- not  \texttt{Occurrences} themselves.}

\subsection{Transfer}\label{subsec:transfer}

\pn \texttt{Transfer} itself now carries the whole mechanism of moving a payload.  \texttt{source}
and \texttt{target} pick up from and drop off to the occurrence's own \texttt{outgoingBuffer} and
\texttt{incomingBuffer}, and three \texttt{step}s, ordered by \texttt{succession}, carry out the
transfer: get the payload, wait \texttt{transferDuration}, then put it.  Both
\texttt{MessageTransfer} (\S\ref{subsec:messagetransfer}) and \texttt{FlowTransfer}
(\S\ref{subsec:flowtransfer}) inherit this unchanged; neither adds anything of its own beyond a
\texttt{doc} comment and mutual \texttt{disjoint}ness.

\ecaption{Transfer}
\begin{lstlisting}
interaction Transfer specializes Performance, BinaryLink {
	end feature source: Occurrence redefines BinaryLink::source {
		feature sourceOutput: Anything[0..*];
	}
	end feature target: Occurrence redefines BinaryLink::target {
		feature targetInput: Anything[0..*];
	}
	feature isInstant: Boolean[1];
	feature transferDuration: DurationValue[1] default 0 [SI::s];
	var feature payload: Anything[1];

	inv { payload istype Physical implies size(source.outgoingTransfersFromSelf) <= 1 }

	succession first getPayloadFromSource then delay;
	succession first delay then putPayloadToTarget;

	step getPayloadFromSource: source.outgoingBuffer.get {
		out payload;
	}
	step delay : TriggerAfter {
		in delay = transferDuration;
		in receiver = this;
	}
	step putPayloadToTarget: target.incomingBuffer.put {
		in payload;
	}
}
\end{lstlisting}

\pn The invariant is a physical-payload restriction, not a temporal one: a \texttt{Physical}
payload cannot be in two places at once, so if it is what is flowing, its source may feed at most
one outgoing transfer.  Only \texttt{Virtual} payloads may be multicast, via
\texttt{outgoingBuffer}'s own non-destructive \texttt{get}.

\subsection{Message Transfer}\label{subsec:messagetransfer}

\pn \texttt{MessageTransfer} and \texttt{FlowTransfer} are \texttt{disjoint} specializations of
\texttt{Transfer}, distinguished by how they are used, not by different features -- both now
inherit \texttt{Transfer}'s complete pickup/delay/dropoff mechanism unchanged, so it is not quite
right to say \texttt{MessageTransfer} alone ``does not specify where the payload is picked up and
dropped off'': neither \texttt{MessageTransfer} nor \texttt{FlowTransfer} adds anything to
\texttt{Transfer}'s own \texttt{sourceOutput}/\texttt{targetInput} at present.
\texttt{MessageTransfer} is the one \texttt{SendPerformance} and \texttt{AcceptPerformance}
(\S\ref{subsec:sendaccept}) use explicitly; \texttt{FlowTransfer} is what a \texttt{flow} connector
is typed by implicitly.

\ecaption{MessageTransfer}
\begin{lstlisting}
	interaction MessageTransfer specializes Transfer {
		doc
		/*
		 * A MessageTransfer is a Transfer that does not specify where the payload is picked
		 * up and dropped off (see FlowTransfer). They are sent by SendPerformances and
		 * accepted by AcceptPerformances.
		 */
	}
\end{lstlisting}

\subsection{Flow Transfer}\label{subsec:flowtransfer}
\ecaption{FlowTransfer}
\begin{lstlisting}
	interaction FlowTransfer specializes Transfer disjoint from MessageTransfer {
		doc
		/*
		 * A FlowTransfer is a Transfer identifying an output feature of the source from which
		 * to pick up a payload and an input feature of the target to which to drop it off. They can
		 * start when the payload is available at the source and move or copy it to the target.
		 */
	}
\end{lstlisting}

\pn \texttt{FlowTransfer}'s own doc comment still describes identifying specific output/input
features of the source/target, and moving or copying a payload -- but the feature content that
would express that (\texttt{isMove}, \texttt{isPush}, and \texttt{sourceOutputLink}/
\texttt{targetInputLink} connectors wiring to specific features) is currently commented out in the
real library.  What is live is exactly \texttt{Transfer}'s own mechanism above: every transfer,
flow or message, picks up from \texttt{source.outgoingBuffer} and drops off to
\texttt{target.incomingBuffer} as a whole, after \texttt{transferDuration}.

\subsection{Send and Accept Performances}\label{subsec:sendaccept}

\pn \texttt{SendPerformance} loads a payload into the sender's \texttt{outgoingBuffer}, then moves
it to the receiver with a \texttt{MessageTransfer}.

\ecaption{SendPerformance}
\begin{lstlisting}
	behavior SendPerformance specializes Performance  {
		in feature payload[1];
		in feature sender: Occurrence[1] default this;
		in feature receiver: Occurrence[1];

		succession loadOutgoingBuffer then sendTransfer;

		step loadOutgoingBuffer : sender.outgoingBuffer.put {
			in value = payload;
		}

		step sendTransfer: MessageTransfer [1] {
			feature redefines payload = SendPerformance::payload;
			end feature redefines source = sender;
			end feature redefines target = receiver;
		}
	}
\end{lstlisting}

\pn \texttt{AcceptPerformance} waits until the receiver's \texttt{incomingBuffer} contains
something conforming to the type of \texttt{payload}, then selects and removes that item --
implementing SysML v2 \S8.4.13.6: the semantics of \texttt{AcceptPerformance} is to select an
incoming transfer whose transferred value conforms to its accepted item parameter, rather than
unconditionally taking whatever is at the head of the buffer.

\ecaption{AcceptPerformance}
\begin{lstlisting}
	behavior AcceptPerformance specializes Performance {
		inout feature payload[1];
		in feature receiver: Occurrence[1] default this;
		feature acceptedTransfer: MessageTransfer[1];
		succession first waitForTransfer then getPayloadFromBuffer;

		step waitForTransfer : GetChangeToFalse {
			in feature d = receiver {
				feature :>> e : BooleanEvaluation[1] {
					step matching : receiver.incomingBuffer.matching {
						in sample = payload;
					}
					out feature redefines result = isEmpty(matching.matched);
				}
			}
		}

		step getPayloadFromBuffer {
			step matching : receiver.incomingBuffer.matching {
				in sample = payload;
			}
			private feature matched : Anything[0..1] = matching.matched;
			step : receiver.incomingBuffer.evict {
				in value = matched;
			}
			binding payload = matched;
		}
	}
\end{lstlisting}

%\comment{Multiple issues:  ItemFlow specializes Connector and Step, so it's a Classifier.  But the checkItemFlowSpecialization constraint requires that ItemFlows specialize Transfers::transfers, which is a \emph{feature}.  The metaclass definition of ItemFlow does not correspond to Transfers::Transfer.}
%
%\begin{lstlisting}
%		metaclass ItemFlow specializes Connector, Step {
%			derived feature itemType : Classifier[0..*];
%			derived feature targetInputFeature : Feature[0..1];
%			derived feature sourceOutputFeature : Feature[0..1];
%			derived feature itemFlowEnd : ItemFlowEnd[0..2] subsets connectorEnd;
%			derived feature itemFeature : ItemFeature[0..1] subsets ownedFeature;
%			derived feature 'interaction' : Interaction[0..*] redefines association, 'behavior';
%		}
%\end{lstlisting}
%
%\begin{lstlisting}
%	interaction Transfer specializes Performance, BinaryLink {
%		end feature source: Occurrence[0..*] redefines BinaryLink::source {
%			feature sourceOutput: Occurrence[0..*];
%		}
%		end feature target: Occurrence[0..*] redefines BinaryLink::target {
%		    feature targetInput: Occurrence[0..*];
%		}
%		readonly feature isInstant: Boolean[1] { }
%		feature item: Anything[1..*] { }
%		feature itemNum: Natural [1] = size(item);
%		private instantNum: Natural[1] = if isInstant? 1 else 0;
%		private binding instant[instantNum] of startShot[0..1] = endShot[0..1] { }
%	}
%\end{lstlisting}

%\comment{ItemFlow::itemType is typed by Classifier; Transfer::item can be Anything.}


